\section{What a linear-time transposer needs}\label{sec:construction}

This section reads the results backwards, as a specification that any oblivious linear-time transposer must meet, and then
describes how close known schedules come to meeting it.

\subsection{The specification}

Suppose $M$ is an oblivious machine with $k$ work tapes that transposes in time $C(L+1)$ at infinitely many sizes. By the
results of this paper and of the series, at all large such sizes:
\begin{enumerate}[leftmargin=2em,label=(\roman*)]
\item \emph{Shape.} At scale $\theta L/4$ each work tape's walk is one of finitely many words, and some tape turns at least
twice with overlapping legs, so it reuses a stretch of tape like a queue (\cref{thm:S}). Along infinitely many sizes the shape
is one fixed tuple of words.
\item \emph{Crossings.} From every set of at least $\alpha L$ outputs, a linear number of vertex-disjoint dependency paths reach
the earlier ends of crossing revisits of excursion more than $\theta L$ (\cref{thm:X}), and the crossing steps of one coarsely monotone leg
of one tape already carry a linear fraction of them; the tape and the leg may depend on the output set and on the size, and are fixed along
infinitely many sizes for the set of all outputs (\cref{thm:S}).
\item \emph{Two long-range heads.} At least two work heads revisit cells after moving a fixed fraction of $L$ away, and for each work head
some set containing a fixed fraction of the outputs reaches long revisits off that head with linear expansion (\cref{thm:H}).
\item \emph{Rows to columns.} For every two sets of rows $A_0,A_1$, every set of steps separating the reads of rows $A_0$ from
the outputs that rows $A_0$ feed into output rows $A_1$ has at least $\abs{A_0}\abs{A_1}b/(\wval+1)$ elements (\cref{lem:R}).
\item \emph{Codes.} After a genie of cost $\eps L$ and bounded depth, the cone relations of the schedule carry an index code of at most
$\eps L$ bits, one of whose fibres is a storage code of rate at least $1-\eps$~\cite[Corollary~O]{BoundedStack2026}; by (i)--(iii), the
relations involved are those of macroscopic queue-like reuse.
\end{enumerate}
Conditions (i)--(iv) are proved here without any coding hypothesis. Condition (v) is the coding question itself: a linear-time
oblivious transposer is a code of rate close to $1$ on the cone relations of a schedule meeting (i)--(iv). Machine-free relations carry
storage codes above rate one half~\cite{StorageCodes2026}, and binary codes of rate approaching one exist on suitable separated
relations~\cite{HighRate2026}.

\subsection{Mechanisms that fail}

Each of the following design patterns is excluded for oblivious linear-time transposition, whatever the machine computes on
its tapes.
\begin{itemize}[leftmargin=1.5em]
\item \emph{Stack-like reuse at macroscopic scale.} A schedule whose long revisits are nested, so that every tape reuses its
storage like a stack at scale $\theta L$, has no macroscopic crossing steps, and is excluded by \cref{thm:X}.
\item \emph{One long-range head.} One unrestricted work head together with heads whose revisits stay below a small fixed fraction
of $L$ is excluded by \cref{thm:H}, whatever the input head does.
\item \emph{Boundedly many macroscopic phases.} Schedules in which the input head and some bank heads are weakly monotone on boundedly many
phases while the other heads move inside sublinear workspaces are excluded by \cref{thm:BM}; with sublinear jitter on the bank heads, and
for adaptive machines, by \cref{thm:J}.
\item \emph{Cheap row-to-column cuts.} A schedule with a set of fewer than $\abs{A_0}\abs{A_1}b/(\wval+1)$ steps separating the reads
of some rows $A_0$ from their outputs in output rows $A_1$ is excluded by \cref{lem:R}; \cref{thm:colouring} gives a convenient
way to exhibit such cuts.
\item \emph{The explicit families of \cref{app:catalogue}}, each excluded by a class theorem of \cref{sec:classes} or a certificate of
\cref{sec:certificates} (in proof sketches). Several of them make macroscopic crossings with linear expansion (numerically, at $d\le128$), so
\cref{thm:X} alone does not exclude them.
\end{itemize}

\subsection{Schedules that meet every condition we could test}\label{sec:cascades}

Five explicit oblivious schedules, which we call \emph{two-sided cascades}, meet every necessary condition of this paper that we could
test. Each loads the input onto two or three work tapes and then runs a cascade over the scales $M_r=L/2^r$, $1\le r<\log_2d$. A forward
stage folds the data at each scale: one half by a transfer at varying relative speed, the other merged in with a cyclic shift of about
$M_r/3$. A backward stage unfolds it by $1$-to-$2$ doublings with an unshifted and a shifted half. A final sweep emits the $L$ output bits.
Every stage at scale $r$ moves the heads over $O(M_r)$ cells a bounded number of times, so the running time is $O(\sum_rM_r)=O(L)$. The five
differ as follows.
\begin{enumerate}[leftmargin=2em]
\item Two work tapes; before the cascade, two symmetric pairs of two-dimensional cyclic shifts of the entry grid are applied to a copy of
the input.
\item Two work tapes; before the cascade, an anti-diagonal reflection and two butterfly shifts, by $2^{m-1}(d-1)b$ and $(d-1)b$ where
$d=2^m$, are applied to a copy.
\item Three work tapes, every cell entered at most ten times; the third tape holds disjoint zones for the folded data at all scales and
moves through them monotonically in each stage.
\item Three work tapes; the forward and backward stages use disjoint zones for all scales, and no zone is overwritten.
\item Two work tapes; the folds are $2$-to-$1$ compressions, and every scale has a comb pass in which one tape makes $2^r$ teeth while the
other waits.
\end{enumerate}
The machines that realise them in our experiments compute other functions; whether some machine with one of these oblivious traces
transposes is \cref{prob:cascades}.

\begin{remark}[Numerical]\label{rem:cascades}
At $d\le128$ and $b=1$ the measured running times of the five schedules lay between $18L$ and $90L$ and stabilised as $d$ grew. On all five,
the minimum set separating a large output set from the earlier ends of crossing revisits of excursion greater than $E$ had exactly as many
elements as the output set, for every tested $E\le L/2$. Their skeletons at scale about $L/6$ were independent of $d$, with a single leg
carrying linear expansion, and their numbers of disjoint expanding sets of steps stayed bounded, as linear time requires. So, as far as
these measurements go, the conditions (i)--(iv) of the specification do not distinguish these schedules from a transposer. A one-pass shear,
excluded by \cref{thm:J}, and a multi-scale stack schedule, excluded by a certificate (in a proof sketch, \cref{app:catalogue}), also make
macroscopic crossings with linear expansion in these measurements, which suggests that (i)--(iv) are not sufficient for transposition.
\end{remark}


\subsection{Why cuts alone cannot finish}

The certificates of this paper are cuts: they bound from below the number of steps that must separate inputs from outputs. A
single cut can only certify storage, and the next proposition shows that for transposition one set of steps serves as the cut
for every pair of row sets at once.

\begin{proposition}[One cut for every pair of row sets]\label{prop:onecut}
\notformalized
Let a run with horizon $T$ read its last input cell at step $t_{\mathrm{in}}$, and put $t^*=t_{\mathrm{in}}+1$. Let $W^*$ consist of
\begin{itemize}[leftmargin=1.5em]
\item every step $u<t^*$ that is the earlier end of a revisit $(u,t)$ with $t\ge t^*$,
\item every output step before $t^*$, and
\item the step $t^*$, if $t^*<T$.
\end{itemize}
Then for every two sets $A_0,A_1$ of rows, $W^*$ meets every dependency path from the output steps of the entries that rows $A_0$
feed into output rows $A_1$ to the steps that read rows $A_0$. The steps of the first kind number at most the number of work cells that
are visited both before $t^*$ and at or after~$t^*$.
\end{proposition}

\begin{informalproof}
Take a dependency path from an output step $t_0$ to an input-reading step $u$. Every edge of the dependency DAG goes from a step to
a strictly earlier step: the state edge $a\to a-1$, or the revisit edge $a\to a'$ for a revisit $(a',a)$. If $t_0<t^*$, then $t_0$ is
an output step before $t^*$ and lies in $W^*$. Otherwise $t_0\ge t^*>t_{\mathrm{in}}\ge u$, so the path has an edge $a\to a'$ with
$a\ge t^*>a'$. If it is a state edge, then $a'=a-1$, so $a=t^*$, and $t^*\in W^*$. If it is a revisit edge, then $(a',a)$ is a revisit
with $a'<t^*\le a$, so $a'\in W^*$. In each case the path meets $W^*$; this argument does not use $A_0$ or $A_1$. A step of the first kind
is the earlier end of a revisit of some tape that spans $t^*$; on each tape, each cell is the cell of at most one such revisit, namely
the one from its last visit before $t^*$ to its next visit.
\end{informalproof}

For a transposer, \cref{lem:R} with $A_0=A_1$ the set of all rows gives $\abs{W^*}\ge(1-\eta)L/(\wval+1)$. When few outputs precede the last
input read, this says that the storage live across $t^*$ is linear. In any case \cref{prop:onecut} shows that this one bound already
accounts for every pair of row sets: however finely the rows are split, cut certificates for different splits may coincide, and a lower bound
beyond linear time cannot come from counting disjoint cuts unless something other than the cut property forces the cuts apart. The
colouring criterion of \cref{sec:demand} verifies such a separation on specific schedules (\cref{cor:scale-separated}), but it does not supply
one in general.

Outside such restricted classes, the only coding-proof logarithmic factor known on tapes is of a different kind. For an oblivious transposer
whose run has $r$ monotone passes, Theorem~P of the first paper gives $k\,db\le((k+1)^r-1)(2b+\sigma_M)$, where $\sigma_M$ is the number of
bits of a state and the scanned symbols, so a fixed machine needs $r\ge\log_{k+1}d-O(1)$ passes~\cite[Theorem~5.9]{FirstPaper2026}; for
adaptive machines an extra term appears~\cite[Theorem~5.12]{FirstPaper2026}. Its grouped form~\cite[Theorem~5.8]{FirstPaper2026} reads
$d^2b\le H_{\mathrm{addr}}\bigl((\gamma+1)db+\lceil d/\gamma\rceil\sigma_M\bigr)$ for every $\gamma\ge1$, where $H_{\mathrm{addr}}=((k+1)^r-1)/k$
bounds the number of addresses of the forests in its proof.

\begin{remark}[Pass counting at one scale]\label{rem:onescale}
Proof sketch only: suppose a run splits into $r$ phases in each of which every head stays within distance $S$ of a monotone trend. Adapting
the proof of the grouped bound so that the nodes of its forests carry the windows of $O(S)$ cells around the trends, at the cost of
bookkeeping for the offsets within each phase, plays the role of a finite control of $\sigma\approx\sigma_M+O(kS\log(g+1))$ bits, and choosing
$\gamma\approx\sqrt{\sigma/b}$ in the grouped inequality gives $(k+1)^r\gtrsim\tfrac k2\sqrt{L/\sigma}$. So boundedly many phases with $S=o(L)$
cannot transpose. At $S=\eps L$, however, the bound only asks for $r\ge\frac12\log_{k+1}(1/\eps)-O(1)$ phases. A linear-time run may have up to
about $4kC/\eps$ legs at scale $\eps L$ (\cref{lem:finite}), and schedules whose number of legs at scale $\eps L$ grows like $\log(1/\eps)$, such
as the cascades of \cref{sec:cascades}, meet the bound at every scale.
\end{remark}

So pass counting, like cut counting, cannot exclude linear time one scale at a time. \Cref{prob:multiscale} asks for a pass bound
that counts phases across all scales at once.
